% Finite free-reference rank and point-probe calculation.
\section{Free-reference rank from finite point probes}
\label{sec:ex-lattice-trace-free-rank}

This section proves the finite-dimensional reference-rank facts used in the
trace calibration. The calculation uses only four-dimensional Euclidean
Green-function derivatives, the signed coordinate symmetries, and simplicity
of the compact gauge Lie algebra. All spaces below are real, even though
polynomial coefficients are temporarily complexified to encode real
observables.

\subsection{Low-grade invariant labels}

Write $\mathfrak g$ for the compact simple real gauge Lie algebra, with its
invariant inner product, and $e_0,e_1,e_2,e_3$ for the coordinate basis of
$\mathbb R^4$. A grade-$d$ label is a polynomial in all radial jets, with jet order $r$
having engineering weight $r+2$. In particular, grades four and five may
include linear terms in higher jets; those terms are removed below by gauge
simplicity. After that reduction, their nonlinear terms depend only on the
curvature $F$ and, at grade five, its first derivative $J$. A signed
permutation $s$ acts
on a coordinate tensor by permuting its indices and multiplying by the
product of the corresponding coordinate signs. The raw grade-four space
consists of symmetric two-index arrays of grade-four labels, real under
complex conjugation and covariant under this signed action. The raw grade-five
space consists of one-index arrays of grade-five labels with the same reality
and covariance conditions.

The following elementary color fact is used repeatedly. If a real linear
functional $\ell$ on $\mathfrak g$ is invariant under the adjoint action,
then differentiating $\ell(\operatorname{Ad}_{\exp(tx)}y)=\ell(y)$ at $t=0$
gives $\ell([x,y])=0$. Simplicity and nonabelianity imply
$[\mathfrak g,\mathfrak g]=\mathfrak g$, hence $\ell=0$. For an invariant
real bilinear form $B$, define $T$ by $B(x,y)=\langle Tx,y\rangle$. The
invariance of $B$ and of the positive inner product implies
$T[a,x]=[a,Tx]$ for every $a,x$. Therefore its self-adjoint and
skew-adjoint parts $S=(T+T^*)/2$ and $A=(T-T^*)/2$ both commute with the
adjoint action. Every eigenspace of $S$ is an ideal; the spectral theorem
and simplicity show $S=\lambda I$. Also $-A^2$ is positive self-adjoint
and commutes with the adjoint action, so the same argument gives
$-A^2=cI$ for some $c\ge0$. If $c>0$, then $J=A/\sqrt c$ is a skew-adjoint
complex structure commuting with every $\operatorname{ad}_a$. This implies
$\operatorname{ad}_{Jx}=J\operatorname{ad}_x$. The right side is
self-adjoint, because it is a product of commuting skew-adjoint operators,
whereas the left side is skew-adjoint by invariance of the inner product.
It must vanish; invertibility of $J$ would then make every
$\operatorname{ad}_x$ vanish, contradicting nonabelianity. Thus $c=0$ and
$A=0$, so $T=\lambda I$. Every invariant bilinear form is consequently a
scalar multiple of the inner product. Applying this argument to real and
imaginary parts covers complex-valued coefficients too.
Consequently the color dependence of an invariant quadratic term is just
the inner product. The engineering weights are $2$ for $F$ and $3$ for its first derivative $J$; a jet of order $r$ has weight $r+2$. At grade four or five, a monomial can therefore contain at most two factors. The possible terms are a constant, one linear functional of each jet of weight at most the grade, a symmetric quadratic form in $F$, and (only at grade five) a mixed bilinear form in $(F,J)$.
Every linear jet term vanishes by the preceding color argument, including the higher-derivative linear terms allowed by these two grades. Under simultaneous
reversal of all four coordinate signs, $F$ is even while the output vector
index and $J$ are odd. Hence constants and quadratic $F$ terms vanish in a
covariant vector. Only the mixed term remains in grade five.

For precision, the source represents the radial degree-zero and degree-one
potentials by the following spanning families, for $c\in\mathfrak g$:
\begin{align*}
 (\ell^0_{ij}c)_a(z)&=\tfrac12(\delta_{aj}z_i-\delta_{ai}z_j)c,\\
 (\ell^1_{rkl}c)_a(z)&=\tfrac13(\delta_{al}z_k-\delta_{ak}z_l)z_r c.
\end{align*}
The first family is antisymmetric in $i,j$; the second is antisymmetric in
$k,l$ and satisfies
\[
 \ell^1_{rkl}+\ell^1_{klr}+\ell^1_{lrk}=0.
\]
This identity follows by adding the three displayed quadratic monomials:
each term cancels its opposite. The families span the actual radial jet
spaces: for every degree-zero potential $A_0$ and degree-one potential
$A_1$, respectively,
\[
 A_0=\tfrac12\sum_{i,j}\ell^0_{ij}(\operatorname{curl}_0^{ij}A_0),
 \qquad
 A_1=\tfrac12\sum_{r,k,l}\ell^1_{rkl}(\operatorname{curl}_1^{rkl}A_1).
\]
Indeed, applying the corresponding curl to the right sides gives the
identity on each curl coordinate, and the displayed formulas for the legs
then give equality of the potential coefficients. Thus a bilinear form in
these jet slots is determined by its values on pairs of legs. Signed
covariance and pair antisymmetry follow directly from the same formulas.
Put $f_{ij}=\operatorname{curl}_0^{ij}A_0$ and
$j_{rkl}=\operatorname{curl}_1^{rkl}A_1$; the cycle identity is precisely
the Bianchi relation $j_{rkl}+j_{klr}+j_{lrk}=0$.

For a grade-five vector label, let $M_p$ denote the mixed bilinear term
in its low-grade decomposition. Choose a unit color vector $v$ and define
\[
 A_{p;r;ij;kl}=M_p(\ell^0_{ij}(v),\ell^1_{rkl}(v)).
\]
Reality makes these coefficients real. The adjoint-invariant bilinear-form
calculation shows that this unit color vector records the full color
coefficient. The leg spanning identities show that the coefficients
determine every $M_p$. The cycle relation for $\ell^1$ gives the Bianchi
relation for $A$. Hence $A_{p;r;ij;kl}$ is antisymmetric in $i,j$ and in
$k,l$, and obeys
\[
 A_{p;r;ij;kl}
 =s_p s_r s_i s_j s_k s_l\,
 A_{s(p);s(r);s(i)s(j);s(k)s(l)}
\tag{1}\label{eq:ex-trace-signed-array}
\]
for every signed permutation, where the signs $s_a\in\{\pm1\}$ are those
of the coordinate reflection. The finite signed-coordinate orbit count in
dimension four has the following explicit form: if a coefficient satisfying
(1) is nonzero, then its two antisymmetric pairs have distinct entries, and
either (i) $p=r$ and the two pairs are parallel, or (ii) $p\ne r$ and there
is a third coordinate $u$, distinct from $p,r$, such that the two pairs are
$(p,u),(r,u)$ in either order. Permuting coordinates and reversing an
antisymmetric pair reduces these cases to the four rows
\[
 (p,r;ij;kl)=(0,0;01;01),\ (0,0;12;12),\
 (0,2;01;21),\ (0,2;21;01).
\tag{2}\label{eq:ex-trace-orbit-representatives}
\]
To check the count directly, a reflection in coordinate $a$ multiplies the
coefficient by $-1$ exactly when $a$ belongs to an odd number of the six
indices $p,r,i,j,k,l$; a nonzero entry requires even parity at every
coordinate. If $p=r$, this forces the supports of the two antisymmetric
pairs to agree, giving case (i). If $p\ne r$, parity forces the pairs to
share exactly one coordinate, distinct from both $p,r$, giving case (ii).
Coordinate permutations put the resulting supports in the forms in (2),
and antisymmetry fixes the order signs. This proves the stated orbit count.

\begin{lemma}[Three probes determine every covariant vector label]
\label{lem:ex-trace-vector-probes}
Let $A_{p;r;ij;kl}$ be a real array satisfying the antisymmetries,
Bianchi relation, and signed covariance above. If
\[
 A_{0;0;01;01}=A_{0;2;01;21}=A_{0;2;21;01}=0,
\]
then every coefficient of $A$ vanishes.
\end{lemma}
\begin{proof}
By the orbit count, it suffices to check the four rows in
\eqref{eq:ex-trace-orbit-representatives}. Three are the displayed
hypotheses. For the remaining row, use Bianchi in the differentiated pair
with derivative index $0$ and pair $12$:
\[
 A_{0;0;12;12}+A_{0;1;12;20}+A_{0;2;12;01}=0.
\]
The third term is minus the third displayed probe by reversing the first
pair $12$ to $21$. Interchanging $1$ and $2$ carries the second term to
the same third probe, up to the sign from reversing its first or
differentiated pair. Both terms therefore vanish, so $A_{0;0;12;12}=0$.
Thus all four representatives vanish; signed
covariance and antisymmetry propagate zero to every entry.
\end{proof}

\subsection{The raw spaces have dimensions at most four and three}

For a raw grade-four tensor $Q=(Q_{ab})$, choose a unit color vector $v$.
The four real linear probes are: the real constant coefficient of $Q_{00}$;
and the coefficients of $\langle F_{01},F_{01}\rangle$ and
$\langle F_{12},F_{12}\rangle$ in $Q_{00}$ and
$\langle F_{01},F_{21}\rangle$ in $Q_{02}$. If all four vanish, the
constant vanishes. Polarization in the two $F$ slots extracts the symmetric
quadratic coefficient. For the quadratic coefficient array, signed parity says that if $a=b$,
the two antisymmetric field pairs have the same support. Relative to the
fixed output axis $a$, their common support either contains $a$ or avoids
$a$; coordinate permutations fixing $a$ reduce these to the probes
$Q_{00}$ on $(01,01)$ and $(12,12)$. If $a\ne b$, the two field-pair
supports have symmetric difference $\{a,b\}$, so they share a third
coordinate distinct from $a,b$. Coordinate permutation reduces this to
$Q_{02}$ on $(01,21)$. Pair reversals change only signs; exchanging the
output indices gives the same coefficient by $Q_{ab}=Q_{ba}$. These are
exactly the three quadratic probes listed above.
The invariant-color argument makes vanishing on the chosen unit color vector
sufficient: every color coefficient is its inner product multiple. The
linear term is zero by color simplicity. Thus these four probes are
injective, and
\[
 \dim \mathcal R_4\le4,
\]
where $\mathcal R_4$ denotes the raw symmetric grade-four table space.

For a raw grade-five vector $Q=(Q_p)$, polarize the mixed term by evaluating the grade-zero and grade-one slots at
$(F,J)$, $(F,0)$, $(0,J)$, and $(0,0)$ and taking the alternating sum.
This cancels the constant and pure-degree terms and leaves exactly the
mixed bilinear coefficient. The three probes
are the three entries in Lemma~\ref{lem:ex-trace-vector-probes}. The Bianchi
identity supplies the fourth orbit representative, while the color argument
and the all-coordinate reflection remove the other grades as shown above.
Thus the probe map is injective into $\mathbb R^3$ and
\[
 \dim \mathcal R_5\le3.
\]
This proves the two raw-dimension bounds without assuming a tensor
classification beyond the finite parity count just given.

\subsection{The normalized free stress basis}

Let $\nu=(4\pi^2)^{-1}$ and $d_G=\dim\mathfrak g>0$. The normalized
four-dimensional massless Green function is $G(x)=\nu|x|^{-2}$. At the unit
vector $e_0$, direct differentiation gives
\begin{align}
 \partial_i\partial_jG(e_0)
   &=-2\nu\delta_{ij}+8\nu\delta_{i0}\delta_{j0},
 \label{eq:ex-trace-green-hessian}\\
 \partial_k\partial_i\partial_jG(e_0)
   &=8\nu(\delta_{ij}\delta_{k0}+\delta_{ki}\delta_{j0}
                 +\delta_{i0}\delta_{kj})
       -48\nu\delta_{i0}\delta_{j0}\delta_{k0}.
 \label{eq:ex-trace-green-third}
\end{align}
For example, differentiating $G$ once gives
$\partial_iG(x)=-2\nu x_i|x|^{-4}$. Differentiating again gives
\[ -2\nu\delta_{ij}|x|^{-4}+8\nu x_ix_j|x|^{-6}. \]
One more derivative and
$x=e_0$ give the displayed formulas.

For antisymmetric index pairs define the free curl kernel
\[
 C_{ab;km}(x)=\delta_{kb}\partial_a\partial_mG(x)
 -\delta_{mb}\partial_a\partial_kG(x)
 -\delta_{ka}\partial_b\partial_mG(x)
 +\delta_{ma}\partial_b\partial_kG(x),
\]
and the two-contraction chain
\[
 \mathcal C_{ab,cd;ij,kl}(x)
 =C_{ab;ij}(x)C_{cd;kl}(x)
  +C_{ab;kl}(x)C_{cd;ij}(x).
\]
The Wick pairing of two separated point insertions with coefficient vector
$w=(w_0,w_1,w_2,w_3)$ is $d_G$ times
\[
 \mathcal D_w^{ab}(H)=
 \begin{cases}
 w_1\sum_i H_{ai,ai}+w_2\sum_{i,j}H_{ij,ij},&a=b,\\
 w_3\sum_i H_{ai,bi},&a\ne b,
 \end{cases}
\]
For fixed field pairs $(ij,kl)$ set
$H_{pqrs}(x)=C_{pq;ij}(x)C_{rs;kl}(x)+C_{pq;kl}(x)C_{rs;ij}(x)$.
The probe is $S_w(ij,kl;ab;x)=d_G\mathcal D_w^{ab}(H(x))$; the vacuum
contribution to a one-point $00$ entry is $w_0$. Substitution of
\eqref{eq:ex-trace-green-hessian} into the curl factors gives
\begin{align}
 S_w(01,01;00;e_0)&=d_G\nu^2(32w_1+64w_2),\\
 S_w(12,12;00;e_0)&=d_G\nu^2(64w_2),\\
 S_w(02,12;01;e_0)&=d_G\nu^2(-16w_3).
 \label{eq:ex-trace-vacuum-probes}
\end{align}
After factoring out $d_G\nu^2$, the coefficient rows for these three
probes, in columns $(w_1,w_2,w_3)$, are
\[
 (32,64,0),\qquad(0,64,0),\qquad(0,0,-16).
\]
These finite sums are evaluated by inserting the displayed Kronecker-delta
Hessian into $C$ and then into $H$ and $\mathcal D_w$. The second equation
recovers $w_2$, the first then recovers $w_1$, and the third recovers $w_3$,
since $d_G\nu^2>0$. Together with the one-point vacuum coefficient $w_0$,
these four probes prove injectivity of the free stress coefficient map into
the free stress table. The raw bound $\dim\mathcal R_4\le4$ and the
surjective averaging projection from raw arrays to the free table show that
the free table has dimension at most four. Thus the coefficient map is an
isomorphism.

\begin{lemma}[Free stress coefficient map is an isomorphism]
\label{lem:ex-trace-free-stress-basis}
The four coefficient tensors consisting of the vacuum scalar, the two
independent diagonal curvature contractions, and the off-diagonal shear
contraction form a basis of the free symmetric grade-four stress table.
\end{lemma}
\begin{proof}
The four probes above prove injectivity from a four-dimensional space; the raw
upper bound proves equality of dimensions, hence bijectivity.
\end{proof}

\subsection{The two conserved free directions}

Let $\mathcal T_4$ be the free symmetric grade-four stress table and
$K_0=\ker(\operatorname{skim}:\mathcal T_4\to\mathbb R)$. Let
$\mathcal V_5$ be the grade-five vector table after imposing the free
Maxwell equation, and let $D_0:K_0\to\mathcal V_5$ be the free
divergence restricted to this zero-skim subspace. Write $\pi$ for the
quotient projection from the raw vector table onto $\mathcal V_5$. The projection is onto by the equivariant
section obtained by averaging any real section over the finite signed
permutation group. The equation-of-motion array is
\[
 E_b=\sum_{i,a}\langle F_{bi},\partial_aF_{ai}\rangle.
\]
It is a nonzero raw covariant vector: take a nonzero color vector $v$, set
$F_{01}=v$ and $\partial_0F_{01}=v$ at one jet, with all other independent
jet coordinates zero. The defining sum then evaluates its $b=0$ component
to $\langle v,v\rangle>0$. The Maxwell equation makes $E$ vanish after the
quotient. Thus the kernel of $\pi$ is nonzero. Rank--nullity and
$\dim\mathcal R_5\le3$ give $\dim\mathcal V_5\le2$.

Let $\phi_1,\phi_2$ be the two linear functionals on $\mathcal V_5$
obtained by the point-divergence probes $(01,01)$ and $(12,12)$ at $e_0$.
For $w$ as above, let $S_w(ij,kl;ab;x)$ denote the response with field-pair
insertion $(ij,kl)$, output component $(a,b)$, and separation $x$. Differentiating
the chain product gives
\[
 \partial_u\mathcal C
 = (\partial_u C_{ab;ij})C_{cd;kl}
  +C_{ab;ij}(\partial_u C_{cd;kl})
  +(\partial_u C_{ab;kl})C_{cd;ij}
  +C_{ab;kl}(\partial_u C_{cd;ij}).
\]
Inserting \eqref{eq:ex-trace-green-hessian}--
\eqref{eq:ex-trace-green-third} and summing $u=0,1,2,3$ gives
\begin{align}
 \sum_u\partial_uS_w(01,01;u0;e_0)
   &=d_G\nu^2(-256w_1-512w_2+128w_3),\\
 \sum_u\partial_uS_w(12,12;u0;e_0)
   &=d_G\nu^2(-512w_2-128w_3).
 \label{eq:ex-trace-divergence-probes}
\end{align}
The two coefficient rows in $(w_1,w_2,w_3)$ are linearly independent.
Consequently $\phi_1\circ D_0$ and $\phi_2\circ D_0$ are independent,
so $\phi_1,\phi_2$ are independent on $\mathcal V_5$. Together with the
upper bound $\dim\mathcal V_5\le2$, this gives $\dim\mathcal V_5=2$,
and $\Phi=(\phi_1,\phi_2):\mathcal V_5\to\mathbb R^2$ is an
isomorphism. The two probe equations have common kernel
$\mathbb R(1,-1/4,1)$ on $(w_1,w_2,w_3)$: the second gives
$w_2=-w_3/4$, and then the first gives $w_1=w_3$. The scalar coefficient
$w_0$ is removed by zero skim. Thus the kernel of $(\operatorname{skim},
\Phi\circ D_0)$ is exactly the Weyl line $(0,1,-1/4,1)$. Since $\Phi$
is an isomorphism, this also proves the Weyl tensor has zero divergence.
The zero-skim domain has dimension three, and the restricted divergence has
one-dimensional kernel, so it has rank two and is onto the two-dimensional
target. Its kernel is precisely the Weyl line.

\begin{proposition}[Free rank and closure dimensions]
\label{prop:ex-trace-free-closure-rank}
The free symmetric grade-four stress table has dimension four. The
Maxwell-reduced grade-five vector table has dimension two. The free
zero-skim divergence map is onto that vector table and its kernel is the
one-dimensional Weyl line.
\end{proposition}
\begin{proof}
The first assertion is Lemma~\ref{lem:ex-trace-free-stress-basis}. The
previous paragraph proves the second and third assertions. The scalar skim
is exactly the real one-point average of the $00$ entry; it is not the tensor
trace. The statement uses only the explicit probe formulas and the
finite-dimensional quotient argument just given.
\end{proof}
